\documentclass[12pt]{article}
\usepackage{fullpage}
\usepackage{amsmath,amssymb,amsthm}
\usepackage[hidelinks]{hyperref}

\newtheorem{theorem}{Theorem}
\newtheorem{lemma}[theorem]{Lemma}
\newtheorem{proposition}[theorem]{Proposition}
\newtheorem{definition}[theorem]{Definition}
\theoremstyle{remark}
\newtheorem{remark}[theorem]{Remark}

\newcommand{\C}{\mathbb C}
\newcommand{\R}{\mathbb R}
\newcommand{\eps}{\varepsilon}
\newcommand{\id}{\mathrm{id}}
\newcommand{\Hom}{\operatorname{Hom}}
\newcommand{\supp}{\operatorname{supp}}
\newcommand{\Symp}{\operatorname{Symp}}
\newcommand{\cT}{\mathcal T}

\title{Problem 4 -- ProofCouncil}
\author{}
\date{}

\begin{document}
\maketitle

\begin{abstract}
Let $(M,\omega)$ be a closed symplectic manifold and $D\subset M$ a
symplectic simple normal crossings divisor. After replacing $\omega$ by
a cohomologous symplectic form joined to it by a path along which $D$
stays simple normal crossings, we construct a tame almost complex
structure $J$ such that every point of $M$ has smooth complex
coordinates $(z_1,\dots,z_n)$ in which $z_1,\dots,z_r$ are
$J$-holomorphic and $D=\{z_1\cdots z_r=0\}$. The proof combines the
regularization theorem of Farajzadeh Tehrani, McLean and Zinger with
auxiliary approximately holomorphic divisors arranged in a forest of
binary trees, whose recursive frames give exact Laurent-monomial
transitions between nested systems of local defining functions, and
with a convexity property of tame structures that patches local
solutions exactly.
\end{abstract}

\section{Introduction}

Throughout, $(M,\omega)$ is a closed symplectic manifold of real
dimension $2n$. A submanifold $V\subset M$ is \emph{symplectic} if
$\omega|_V$ is nondegenerate; symplectic submanifolds, and $M$ itself,
are oriented by their symplectic forms.

\subsection*{Conventions}

\begin{definition}[positive symplectic SNC collection]\label{def:snc}
A finite collection $\{D_i\}_{i\in S}$ of closed embedded symplectic
submanifolds of $M$ of real codimension two is a \emph{positive
symplectic simple normal crossings (SNC) collection} if for every
$\Lambda\subset S$ the \emph{closed intersection}
$D_\Lambda=\bigcap_{i\in\Lambda}D_i$ ($D_\varnothing=M$) has the
following properties: at each point the tangent spaces of the $D_i$,
$i\in\Lambda$, meet transversally, so that $D_\Lambda$ is a (possibly
empty) closed submanifold of codimension $2|\Lambda|$; $D_\Lambda$ is
symplectic; and the orientation of $D_\Lambda$ as a transverse
intersection of the oriented
submanifolds $D_i$ inside the oriented manifold $M$ agrees with its
symplectic orientation. The union $D=\bigcup_iD_i$ is then a
\emph{symplectic SNC divisor} with \emph{branches} $D_i$.
\end{definition}

In the terminology of \cite{FMZ}, $\{D_i\}$ is an ``SC symplectic
divisor'' and $\omega$ belongs to the space $\Symp^+(M,\{D_i\})$ of
symplectic forms positive for it. We work with collections rather than
with the union $D$ because the proof adjoins auxiliary branches. Since
$D_\Lambda$ has codimension $2|\Lambda|$, at most $n$ branches pass
through any point; this elementary fact is used repeatedly.

An almost complex structure $J$ on $M$ is \emph{$\omega$-tame} if
$\omega(v,Jv)>0$ for all $v\neq0$, and \emph{$\omega$-compatible} if in
addition $\omega(Ju,Jv)=\omega(u,v)$. A smooth function $f\colon U\to\C$
on an open set $U\subset M$ is \emph{$J$-holomorphic} if
$df\circ J=i\,df$, equivalently if the complex-linear extension of $df$
vanishes on the $(0,1)$-bundle $T^{0,1}_JM=\{v+iJv\}$. Products, and
quotients by nonvanishing functions, of $J$-holomorphic functions are
$J$-holomorphic. A \emph{smooth complex coordinate system} is a smooth
real coordinate system with values in $\C^n$; it need not be a
holomorphic chart for $J$.

\begin{theorem}\label{thm:main}
Let $(M,\omega)$ be a closed symplectic manifold and $D=\bigcup_{i\in S}D_i$
a symplectic SNC divisor in the sense of Definition~\ref{def:snc}. There
are a smooth path of symplectic forms $\omega_t$, $0\le t\le1$, and an
$\omega_1$-tame almost complex structure $J$ such that $\omega_0=\omega$,
$[\omega_t]=[\omega]\in H^2(M;\R)$ for all $t$, $\{D_i\}$ is a positive
symplectic SNC collection for every $\omega_t$, and every point of $M$
has a neighbourhood with smooth complex coordinates $(z_1,\dots,z_n)$
in which, for some $0\le r\le n$, the functions $z_1,\dots,z_r$ are
$J$-holomorphic and $D=\{z_1\cdots z_r=0\}$.
\end{theorem}

\subsection*{What is being asked}

It is comparatively easy, after the permitted regularization, to find
a tame $J$ for which every branch $D_i$ is a $J$-holomorphic
submanifold, $J(TD_i)=TD_i$ (Lemma~\ref{lem:reg} below; for almost
complex structures adapted to a divisor see also \cite[Section 3]{FZ}).
Theorem~\ref{thm:main} asks for more: a scalar function $f$ with
$df\circ J=i\,df$ on an open set and $D_i=\{f=0\}$ there.

Here is the difficulty. If $D_i$ is the zero set of a transverse
section $s_i$ of a line bundle $L_i$, a local defining function is
$f=s_i/F$ for a local nonvanishing frame $F$. Two frames on
overlapping sets differ by a nonvanishing function which must itself be
$J$-holomorphic if both quotients are; unrelated local trivializations
have arbitrary transition functions, so no single $J$ serves them all.
The central device of the proof is a global system of frames whose
transition functions are exact Laurent monomials in other defining
functions of the same kind.

\subsection*{Strategy}

\emph{Regularization} (Section~\ref{sec:inputs}): by
\cite[Theorem 2.13]{FMZ}, a cohomologous deformation of $\omega$ makes
the normal bundles of the branches symplectically orthogonal along
their intersections; a compatible almost complex structure preserving
every branch follows.
\emph{Auxiliary divisors} (Section~\ref{sec:inputs}): given a line
bundle $L$ and a positive SNC collection $E$, approximately holomorphic
sections \cite{Donaldson1996,Auroux1997,Mohsen2019} produce a line
bundle $H$ and sections $u$ of $H$, $v$ of $L\otimes H$ whose zero sets,
together with $E$, form a positive SNC collection (Lemma~\ref{lem:aux}).
\emph{The forest} (Section~\ref{sec:forest}): over each branch $D_i$ we
grow a full binary tree of depth $n$; a node $a$ carries $(L_a,s_a)$ and
its children carry $L_{a-}=H_a$, $L_{a+}=L_a\otimes H_a$, so that
$L_a\cong L_{a+}\otimes L_{a-}^{-1}$. Where $s_{a\pm}$ do not vanish,
$s_a/(s_{a+}\otimes s_{a-}^{-1})$ is a defining function of
$D_a=s_a^{-1}(0)$; where $s_{a+}$ or $s_{a-}$ is small, that section is
in turn replaced by the frame built from its own children, and so on.
Since at most $n$ divisors pass
through a point, depth $n$ suffices.
\emph{Multiscale cover and exact transitions} (Section~\ref{sec:frames}):
open sets $U_A$, indexed by the patterns $A$ of ``small'' sections,
cover $M$, and on $U_A$ the recursion gives functions $f^A_a$ with zero
sets $D_a\cap U_A$; for $A\subset B$ the $f^A_a$ are Laurent monomials,
with coefficient one, in the $f^B_c$ (Lemma~\ref{lem:monomial}).
\emph{Estimates and local structures} (Section~\ref{sec:estimates}):
with thresholds $\eps_0\ll\dots\ll\eps_n$ attached to the depths, the
differentials $df^B_a$, multiplied by their nonvanishing frames, are
uniformly small perturbations of the nearly complex-linear derivatives
of the sections $s_a$, and a pointwise construction gives, near each
point, a tame structure for which all functions of a maximal pattern
are holomorphic.
\emph{Patching} (Sections~\ref{sec:linear}, \ref{sec:patching}): tame
structures are parametrized by strict contractions (``graph maps'')
relative to a compatible reference structure, and holomorphicity of a fixed
function is an \emph{affine} condition on the graph map, so averaging
graph maps with a partition of unity preserves every common constraint
exactly; Lemma~\ref{lem:monomial} makes the constraints common. The
coordinates are read off in Section~\ref{sec:end}.

\emph{Inputs and credit.}
Besides \cite[Theorem 2.13]{FMZ} and the approximately holomorphic
techniques of \cite{Donaldson1996,Auroux1997,Auroux2001,Mohsen2019},
we use the symplectic polar decomposition \cite[Theorem 2.16]{Wendl}
and S\'evennec's description of tame structures as graphs of strict
contractions \cite[Theorem 2.37]{Wendl}; almost complex structures
preserving auxiliary approximately holomorphic hypersurfaces appear in
\cite[Section 8]{CM}. The remaining ingredients, namely the forest of
auxiliary divisors, the recursive frames with monomial transitions, the
multiscale estimates and the patching with affine constraints, are
developed here.

\section{The symplectic inputs}\label{sec:inputs}

In this section $(X,\sigma)$ is a closed symplectic manifold of
dimension $2n$ and $E=\{E_i\}_{i\in S}$ a positive symplectic SNC
collection in $X$ (Definition~\ref{def:snc}), with closed intersections
$E_\Lambda$, $\Lambda\subset S$. The two lemmas below are consequences
of published results.

\subsection{Regularization}

Farajzadeh Tehrani, McLean and Zinger \cite{FMZ} associate with a
symplectic SNC collection the notion of a \emph{regularization}: a
compatible system of identifications of neighbourhoods of the strata
$E_\Lambda$ with neighbourhoods of the zero sections of their normal
bundles, modelled on split symplectic forms. We use only two
consequences: a regularization makes, along each $E_\Lambda$, the
symplectic normal bundles $N_i=(TE_i)^{\sigma}|_{E_\Lambda}$,
$i\in\Lambda$ (where $V^\sigma$ denotes the $\sigma$-orthogonal
complement of $V$), pairwise $\sigma$-orthogonal, so that
\begin{equation}\label{eq:split}
 TX|_{E_\Lambda}=TE_\Lambda\oplus\bigoplus_{i\in\Lambda}N_i
\end{equation}
is a $\sigma$-orthogonal splitting; and \cite[Theorem 2.13]{FMZ}
(theorem numbers refer to the arXiv version v3), read with its
parameter space a point and without relative conditions, provides a
path of symplectic forms $\sigma_t=\sigma+d\tau_t$, $\tau_0=0$, for which
$E$ remains a positive SNC collection and whose endpoint admits a
regularization.

\begin{lemma}[Regularization]\label{lem:reg}
There is a smooth path $\sigma_t=\sigma+d\tau_t$, $0\le t\le1$, of
symplectic forms with $\tau_0=0$, such that $E$ is a positive symplectic
SNC collection for every $\sigma_t$ and such that for $\sigma_1$ the
splittings \eqref{eq:split} (with $N_i=(TE_i)^{\sigma_1}$) are
$\sigma_1$-orthogonal. Moreover $\sigma_1$ admits a compatible almost
complex structure $K$ with $K(TE_i)=TE_i$ for every $i$.
\end{lemma}

\begin{proof}
The path is \cite[Theorem 2.13]{FMZ}, as just explained. For the
compatible structure, recall the symplectic polar decomposition
\cite[proof of Theorem 2.16]{Wendl}: a Riemannian metric $g$ determines
a $g$-skew endomorphism $A$ by $\sigma_1(u,v)=g(Au,v)$, and
$J_g=A(A^*A)^{-1/2}$ is $\sigma_1$-compatible. If a subspace
$W\subset T_xX$ satisfies $A(W)\subset W$, then also
$A^*(W)=A(W)\subset W$, hence $(A^*A)^{-1/2}$, a limit of polynomials in
$A^*A$, preserves $W$, and so does $J_g$. For $W=T_xE_i$ one has
$A(W)\subset W$ exactly
when $g(Au,v)=\sigma_1(u,v)=0$ for all $u\in W$ and $v\in W^{\perp_g}$,
that is, when the $g$-orthogonal complement of $T_xE_i$ is its
$\sigma_1$-orthogonal complement $N_i(x)$. So it suffices to find a
metric $g$ such that
\begin{equation}\label{eq:metric-cond}
 T_xE_i\perp_g N_i(x)\qquad\text{for all }i\text{ and all }x\in E_i.
\end{equation}
For fixed $x$ and $i$ this is a linear condition on $g$ (orthogonality
of two fixed complementary subspaces), so the metrics on $T_xX$
satisfying \eqref{eq:metric-cond} for all $i$ with $x\in E_i$ form a
convex set. At a point $x\in E_\Lambda$ lying on no deeper stratum this
set consists exactly of the metrics making \eqref{eq:split} orthogonal
at $x$, because
$T_xE_i=T_xE_\Lambda\oplus\bigoplus_{j\in\Lambda\setminus\{i\}}N_j(x)$;
and a metric orthogonalizing the finer splitting of a deeper stratum
$E_{\Lambda''}\ni x$, $\Lambda''\supset\Lambda$, also orthogonalizes
\eqref{eq:split}. Since \eqref{eq:metric-cond} concerns only points of
$\bigcup_iE_i$, define $g$ on $TX|_{E_\Lambda}$ by descending induction
on $|\Lambda|$: the metric already defined over the deeper strata
contained in $E_\Lambda$ orthogonalizes \eqref{eq:split} there; extend
it to a neighbourhood in $E_\Lambda$ summand by summand, and combine it
with local metrics orthogonalizing \eqref{eq:split} by a partition of
unity on $E_\Lambda$, using convexity. Finally extend $g$ to $X$; then
$K=J_g$ is the required structure.
\end{proof}

\subsection{Auxiliary divisors}

We recall the framework of approximately holomorphic sections. Suppose
$[\sigma/(2\pi)]$ is integral, let $A\to X$ be a Hermitian line bundle
with unitary connection of curvature $-i\sigma$, let $K$ be
$\sigma$-compatible and $g=\sigma(\cdot,K\cdot)$. Sections of $A^k$, or
of $F\otimes A^k$ for a fixed Hermitian vector bundle $F$, are measured
in the rescaled metric $g_k=kg$, in which the curvature of $A^k$ has
unit size. A family $s_k$ is \emph{approximately holomorphic} if $s_k$
and its first few covariant derivatives are bounded in $g_k$-norm while
the $K$-antilinear part $\bar\partial_Ks_k$ of $\nabla s_k$ and its
derivatives are $O(k^{-1/2})$; it is
\emph{$\eta$-transverse to zero} along a compact submanifold
$Y\subset X$ if at every $y\in Y$ with $|s_k(y)|<\eta$ the derivative
$\nabla s_k|_{T_yY}$ is surjective with a right inverse of $g_k$-norm
at most $\eta^{-1}$. Donaldson \cite{Donaldson1996} constructed, for
all large $k$, approximately holomorphic sections of $A^k$ that are
$\eta$-transverse along $Y=X$ with $\eta>0$ independent of $k$; Auroux
\cite[Theorem 2 and Corollary 1]{Auroux1997} extended this to
$F\otimes A^k$; Mohsen's relative theorem \cite[p.~657]{Mohsen2019}
achieves $\eta$-transversality of the restriction to an arbitrary
compact submanifold $Y$, not necessarily $K$-invariant, and its
twisted version \cite[p.~662]{Mohsen2019} allows a fixed $F$. The
sections are obtained by perturbations of controlled size, so finitely
many such conditions can be imposed successively, each later
perturbation being smaller than the transversality bounds already
secured (Donaldson's globalization principle, see
\cite[Proposition 4.1]{Auroux2001}).

Transversality produces symplectic submanifolds because of the
following linear-algebra fact, due to Donaldson \cite{Donaldson1996}:
if at a zero $y$ of $s$ the $K$-complex-linear part of $\nabla s$ is
surjective with right inverse of norm at most $2\eta^{-1}$ and the
antilinear part is small compared with $\eta$ (both hold when $\nabla s$
is $\eta$-transverse and the antilinear part is at most $\eta/2$), then
$\ker\nabla s(y)$ is the graph of a small map over a complex subspace,
hence a symplectic subspace whose orientation agrees with the complex
orientation of that subspace. The same holds for the restriction to a
$K$-invariant submanifold $Y$, within $(T_yY,K)$. Since complex
orientations are compatible with transverse intersection, all
intersections obtained in this way are positively oriented in the
sense of Definition~\ref{def:snc}.

\begin{lemma}[Auxiliary divisors]\label{lem:aux}
Let $(X,\sigma)$ be closed, $E$ a finite positive symplectic SNC
collection in $X$, and $L\to X$ a smooth complex line bundle. After a
deformation of $\sigma$ as in Lemma~\ref{lem:reg} there are a complex
line bundle $H$ and smooth sections $u\in\Gamma(H)$,
$v\in\Gamma(L\otimes H)$, transverse to the zero section, such that
$E\cup\{u^{-1}(0),v^{-1}(0)\}$
is a positive symplectic SNC collection. Moreover the orientations
induced on $u^{-1}(0)$ and $v^{-1}(0)$ by the sections $u$, $v$ (through
the complex orientations of their bundles) are the symplectic ones.
\end{lemma}

\begin{proof}
Apply Lemma~\ref{lem:reg}; denote the resulting form again by $\sigma$
and let $K$ be the compatible structure preserving every $E_i$, hence
every closed intersection $E_\Lambda$; put $g=\sigma(\cdot,K\cdot)$.
We work with the compact
$E_\Lambda$ rather than with the open strata because the transversality
theorems require compactness.

\emph{Step 1: the integral case.} Suppose $[\sigma/(2\pi)]$ is
integral, let $A$ have curvature $-i\sigma$ and put $H=A^k$. The pair
$(u,v)$ is a section of $(\underline\C\oplus L)\otimes A^k$. For every
nonempty $E_\Lambda$, including $E_\varnothing=X$, impose
$\eta$-transversality to zero along $E_\Lambda$ for each of $u$, $v$ and
$(u,v)$, using the relative twisted theorem of \cite{Mohsen2019}. These
are finitely many conditions, achieved successively as explained above,
with a common $\eta>0$ for all large $k$. Transversality of
$u|_{E_\Lambda}$ gives $u^{-1}(0)\pitchfork E_\Lambda$; transversality
of $(u,v)|_{E_\Lambda}$ gives $u^{-1}(0)\pitchfork v^{-1}(0)$ along
$E_\Lambda$ and $u^{-1}(0)\cap v^{-1}(0)\pitchfork E_\Lambda$. Where the
expected dimension of an intersection is negative, $\eta$-transversality
means $|s|\ge\eta$ on $E_\Lambda$, so the intersection is empty. Since
$E_\Lambda$ is $K$-invariant, $\nabla(s|_{E_\Lambda})$ splits into
$K$-linear and $K$-antilinear parts, and the antilinear part is
$O(k^{-1/2})$ by approximate holomorphicity;
Donaldson's linear-algebra fact therefore shows that every new
intersection is symplectic and positively oriented, so
$E\cup\{u^{-1}(0),v^{-1}(0)\}$ is a positive SNC collection. The same
fact gives the orientation statement: deforming $\nabla s$ linearly to
its complex-linear part through surjective maps, the orientation
induced by $s$ on its zero set is the complex orientation of the
nearby complex kernel, which is the symplectic orientation.

\emph{Step 2: removing the integrality assumption.} Call a symplectic
form $\sigma'$ \emph{rational} if $[\sigma'/(2\pi)]\in H^2(X;\mathbb Q)$.
In general $\sigma$ is not rational, and we follow the
device of \cite[Section 1]{Auroux2001b}: a nearby rational form is used
only to construct the polarizing bundle, while symplecticity is checked
for $\sigma$ itself. Fix closed $2$-forms $\alpha_j$ representing a
basis of $H^2(X;\R)$ and consider the forms
$\sigma'=\sigma+\sum_jc_j\alpha_j$ with $c$ in a small closed ball; all
are symplectic with uniformly bounded derivatives, and rational forms
are dense among them. Polar decomposition with respect to the fixed
metric $g=\sigma(\cdot,K\cdot)$ gives
$\sigma'$-compatible structures $K'$ depending smoothly on $\sigma'$,
with $K'=K$ at $\sigma'=\sigma$. For a rational $\sigma'$ choose a
positive integer $q$ with $q[\sigma'/(2\pi)]$ integral, let $A$ have
curvature $-iq\sigma'$ and again set $H=A^k$. Since the curvature of
$A^k$ is $-iN\sigma'$ with $N=qk$, the effective large parameter is
$N$ and the effective metric is $g_N=Ng'$, $g'=\sigma'(\,\cdot,K'\cdot\,)$.

Step 1, carried out for $(\sigma',K')$ and the parameter $N$, provides
a transversality constant $\eta$ for the finitely many restriction
conditions. The point requiring attention is that $\eta$ must not
depend on the denominator $q$, because the rational approximation will
be chosen in terms of $\eta$. Uniformity in $q$ is obtained from the
quantitative construction, not merely from qualitative openness; it is
a deduction from the structure of the cited estimates, which we
indicate. Mohsen's linear restriction estimate has constants depending
only on dimensions and radii
\cite[Theorem 6.5 and Corollary 6.6]{Mohsen2019}. In local unitary
gauges at length scale $N^{-1/2}$ the curvature $-iN\sigma'$ is
normalized by $N$, so the geometric bounds required by the construction
are uniform over the compact family of $\sigma'$, and the fixed twist
$L$ contributes only lower-order terms. The local perturbation and
globalization estimates
\cite[Proposition 8.2 and Sections 8.3--8.6]{Mohsen2019} use these
bounds together with counts of separated sets in the fixed compact
$E_\Lambda$. They therefore give a common bound $\eta>0$, independent
of $q$, for all the finitely many conditions when $N$ is large.
%% FLAG: The independence of eta from q is asserted in the baseline on the
%% basis of the dependence of the constants in Mohsen's estimates on
%% normalized local data only. This reading of the cited proofs has not been
%% independently verified in this editorial pass; the attribution review
%% also identified it as the one deduction not stated verbatim in the sources.

The relative theorem does not require $K'$ to preserve $E_\Lambda$, but
$K$ does preserve $E_\Lambda$. For $s=u,v,(u,v)$ the derivative along
$E_\Lambda$ therefore satisfies, in the rescaled metrics,
\[
 \bigl\|(\nabla s|_{TE_\Lambda})^{0,1}_K\bigr\|_{g_N}
 \le C\lVert K'-K\rVert_{C^0}+CN^{-1/2},
\]
the first term coming from the difference between the $K'$- and
$K$-antilinear parts and the second from approximate
$K'$-holomorphicity. Choose the rational approximation $\sigma'$ so
close to $\sigma$ that the first term is small compared with $\eta$,
then take $k$ large to control the second. Donaldson's linear-algebra
fact now applies with respect to $K$, which is $\sigma$-compatible: all
new intersections are $\sigma$-symplectic and positively oriented, not
merely $\sigma'$-symplectic.

The rational form is construction data only and is not inserted into
the deformation path. The sole change of $\sigma$ in this lemma is the
initial regularization deformation of Lemma~\ref{lem:reg}.
\end{proof}

\section{Two linear-algebra facts about tame structures}\label{sec:linear}

For an almost complex structure $J$ on a symplectic manifold $(X,\sigma)$
write $T^{1,0}_JX$ and $T^{0,1}_JX$ for
the $\pm i$ eigenbundles of $J$ on $TX\otimes\C$. The first fact is
S\'evennec's parametrization of tame structures by strict contractions
\cite[Theorem 2.37]{Wendl}, together with the observation that
holomorphicity constraints are affine in this parametrization.

\begin{lemma}[Convex patching with holomorphic-function constraints]
\label{lem:convex}
Let $K$ be a $\sigma$-compatible almost complex structure on $X$ and
$E^+=T^{1,0}_KX$, $E^-=T^{0,1}_KX$. Every $\sigma$-tame almost complex
structure $J$ has $T^{0,1}_JX=\{\xi+\Phi_J\xi:\xi\in E^-\}$ for a unique
complex-linear bundle map $\Phi_J\colon E^-\to E^+$ (the \emph{graph
map} of $J$) with $\lVert\Phi_J\rVert<1$ pointwise, and conversely every
such $\Phi$ arises from a unique tame $J$. For a smooth function $f$,
the condition that $f$ be $J$-holomorphic is the affine linear condition
\begin{equation}\label{eq:affine}
 df(\xi+\Phi_J\xi)=0\qquad(\xi\in E^-)
\end{equation}
on $\Phi_J$. Consequently, if $J_\alpha$ are $\sigma$-tame structures on
open sets $W_\alpha$, $(\rho_\alpha)$ is a partition of unity
subordinate to $(W_\alpha)$, and $J$ is the tame structure with
$\Phi_J=\sum_\alpha\rho_\alpha\Phi_{J_\alpha}$, then every function that
is $J_\alpha$-holomorphic at $y$ for all $\alpha$ with $\rho_\alpha(y)\ne0$
is $J$-holomorphic at $y$.
\end{lemma}

\begin{proof}
On $TX\otimes\C$ consider the Hermitian form
$\mathcal H(\xi,\eta)=-i\,\sigma(\xi,\overline\eta)$. The summands $E^+$
and $E^-$ are $\mathcal H$-orthogonal, $\mathcal H$ is positive
definite on $E^+$ and negative definite on $E^-$, and we use these
restrictions (with the sign reversed on $E^-$) as Hermitian metrics to
define $\lVert\Phi\rVert$. For $\xi=v+iJv\in T^{0,1}_JX$ one computes
$\mathcal H(\xi,\xi)=-2\sigma(v,Jv)$, so for tame $J$ the bundle
$T^{0,1}_JX$ is $\mathcal H$-negative definite of complex rank $n$.
Such a subspace meets the positive definite $E^+$ trivially, hence is
the graph of a complex-linear map $\Phi_J\colon E^-\to E^+$, and
negativity of $\mathcal H(\xi+\Phi\xi,\xi+\Phi\xi)=-|\xi|^2+|\Phi\xi|^2$
says exactly $\lVert\Phi_J\rVert<1$. Conversely, the graph of a strict
contraction $\Phi$ is negative definite, its complex conjugate is
positive definite, the two are complementary, and they are the
$(0,1)$- and $(1,0)$-bundles of a real almost complex structure $J$;
the computation above, read backwards, shows that $J$ is tame.

Condition \eqref{eq:affine} is the vanishing of $df$ on $T^{0,1}_JX$
and is affine in $\Phi_J$. A convex combination of strict contractions
is a strict contraction, and a convex combination of solutions of the
affine equations \eqref{eq:affine} for a fixed $f$ is again a solution.
\end{proof}

Averaging the endomorphisms $J_\alpha$ themselves would not produce an
almost complex structure; averaging graph maps does, and no derivative
of the partition of unity enters, so the holomorphicity statements are
exact.

\begin{lemma}[Graph construction]\label{lem:graph}
Let $I$ be a $\sigma$-tame almost complex structure, $V\to X$ a complex
vector bundle with a Hermitian metric, and $\lambda$ a $V$-valued
one-form on an open set $W\subset X$, with complex-linear extension to
$TX\otimes\C$ restricting to
\[
 P=\lambda|_{T^{1,0}_IX}\colon T^{1,0}_IX\to V,\qquad
 Q=\lambda|_{T^{0,1}_IX}\colon T^{0,1}_IX\to V .
\]
Suppose $P$ is surjective on $W$ and let $R$ be a smooth right inverse
of $P$, for instance $R=P^*(PP^*)^{-1}$, which has the least norm among
all right inverses of $P$. Put
\begin{equation}\label{eq:graph}
 \mu=-RQ\colon T^{0,1}_IX\to T^{1,0}_IX,\qquad
 E_\mu=\{\xi+\mu\xi:\xi\in T^{0,1}_IX\}.
\end{equation}
Then $\lambda$ vanishes on $E_\mu$. If $X$ is compact and a fixed
metric is used, there is $\mu_0>0$ such that whenever
$\lVert\mu\rVert<\mu_0$ on $W$, the bundle $E_\mu$ is the $(0,1)$-bundle
of a $\sigma$-tame almost complex structure $J_\mu$ on $W$, depending
smoothly on the data. In particular, if $V=\bigoplus_aL_a$ and the
components of $\lambda$ have the form $\lambda_a=F_a\,df_a$ with $F_a$
a nonvanishing section of $L_a$ over $W$ and $f_a$ a smooth function,
then every $f_a$ is $J_\mu$-holomorphic on $W$.
\end{lemma}

\begin{proof}
For $\xi\in T^{0,1}_IX$, $\lambda(\xi+\mu\xi)=Q\xi+P\mu\xi=Q\xi-PRQ\xi=0$.
With $\mathcal H$ as in the previous proof, $T^{0,1}_IX$ is
$\mathcal H$-negative definite because $I$ is tame, and negativity is
an open condition; on a compact manifold there is a uniform $\mu_0$
such that $E_\mu$ is negative definite whenever $\lVert\mu\rVert<\mu_0$.
Then $E_\mu$ and $\overline{E_\mu}$ are complementary and define a
tame almost complex structure $J_\mu$ with $T^{0,1}_{J_\mu}X=E_\mu$,
exactly as in the proof of Lemma~\ref{lem:convex}. Finally,
$\lambda_a=F_a\,df_a$ vanishes on $E_\mu$ and $F_a$ is nonvanishing,
so $df_a$ vanishes on $E_\mu=T^{0,1}_{J_\mu}X$.
\end{proof}

\section{Building the finite forest}\label{sec:forest}

From now on $X=M$ with the given divisor $D=\bigcup_{i\in S}D_i$, where
$S=\{1,\dots,m\}$. If $D=\varnothing$, Theorem~\ref{thm:main} holds
with $\omega_t=\omega$, any tame $J$ and $r=0$, so assume $D\ne\varnothing$.

\subsection*{Divisor bundles of the branches}
For every branch $D_i$ choose a smooth complex line bundle $L_i\to X$
with a section $s_i$ transverse to the zero section whose zero set is
$D_i$, the orientation induced on $D_i$ by $s_i$ (through the complex
orientation of $L_i$) being the symplectic one: pull the oriented
normal bundle of $D_i$ back to a tubular neighbourhood, take the
tautological section there, and glue to the trivial bundle with the
constant section $1$ over $X\setminus D_i$, using that the tautological
section does not vanish on the overlap (compare \cite[Section 2.2]{FMZS}).

\subsection*{The trees}
Make each $i\in S$ the root of a full binary tree of depth $n$. A node
is a word $a=iw$ with $w\in\{+,-\}^{d}$, $0\le d\le n$; its depth is
$d(a)=d$, roots have depth $0$, \emph{leaves} have depth $n$, and a
non-leaf node $a$ has the two children $a+$ and $a-$. Let $\cT$ be the
finite set of all nodes of all $m$ trees. We attach to every node $a$ a
line bundle $L_a$ with a transverse section $s_a$, and to every non-leaf
node a bundle isomorphism $\theta_a$, by induction on depth.
Suppose $L_a$, $s_a$ have been constructed for all nodes of depth
$\le d<n$, with the zero sets $D_a=s_a^{-1}(0)$ of all nodes constructed
so far forming a positive SNC collection for the current symplectic
form. Process the nodes $a$ of depth $d$ one at a time: apply
Lemma~\ref{lem:aux} with $L=L_a$ and $E$ the collection of all
divisors constructed so far (including those produced for earlier nodes
of depth $d$), and set
\[
 L_{a-}=H_a,\quad s_{a-}=u,\qquad
 L_{a+}=L_a\otimes H_a,\quad s_{a+}=v .
\]
Retain the canonical complex bundle isomorphism
\begin{equation}\label{eq:theta}
 \theta_a\colon L_{a+}\otimes L_{a-}^{-1}
 =L_a\otimes H_a\otimes H_a^{-1}\xrightarrow{\ \cong\ }L_a .
\end{equation}
Here and below inverse line bundles and inverses of nonvanishing
sections are taken in the tensor sense: $F^{-1}$ is the section of
$L^{-1}=L^*$ dual to the nonvanishing section $F$ of $L$. Each
application of Lemma~\ref{lem:aux} includes a regularization
deformation of the symplectic form, which preserves all divisors
constructed so far as a positive SNC collection. There are
$m(2^n-1)$ applications. At the end, the finite collection
$\{D_a\}_{a\in\cT}$ is a positive symplectic SNC collection for a
symplectic form cohomologous to $\omega$, obtained from $\omega$ by
concatenating finitely many paths as in Lemma~\ref{lem:reg}.

Apply Lemma~\ref{lem:reg} once more to the full collection
$\{D_a\}_{a\in\cT}$. Denote the resulting symplectic form by $\sigma$
and let $I$ be the $\sigma$-compatible (in particular $\sigma$-tame)
almost complex structure provided by the lemma, which preserves every
$D_a$. Thus $I$ preserves $TD_a$ and induces a complex structure on
each normal line $N_xD_a=T_xX/T_xD_a$.

\subsection*{Normalizing the sections along their zero sets}
Along $D_a$ the section $s_a$ vanishes, so its derivative there is an
intrinsic $L_a$-valued one-form $ds_a\colon TX|_{D_a}\to L_a|_{D_a}$,
independent of any connection, with kernel $TD_a$. We arrange, by
changing the sections but not their bundles or zero sets, that
\begin{equation}\label{eq:normal-linear}
 (ds_a)\circ I=i\,ds_a\quad\text{on }D_a\qquad(a\in\cT),
\end{equation}
that is, $ds_a$ is complex linear along $D_a$. Indeed, at $x\in D_a$
the derivative descends to a real-linear isomorphism
$N_xD_a\to(L_a)_x$ between complex lines, which may be written
$z\mapsto Az+B\bar z$ with $A,B\in\C$. It is orientation preserving:
the orientation of $D_a$ induced by $s_a$ is the symplectic one (by
construction for roots, by Lemma~\ref{lem:aux} for the other nodes;
both orientations are unchanged along the subsequent deformations),
and since $I$ is tame and preserves $TD_a$, the symplectic orientations
of $X$ and $D_a$ are the complex orientations of $(TX,I)$ and
$(TD_a,I)$, so that $N_xD_a$ carries the quotient complex orientation.
Hence $|A|>|B|$, and by compactness $|A|-|B|$ is bounded below on $D_a$. Identify a tubular neighbourhood of $D_a$ with a
neighbourhood of the zero section of $ND_a$ and identify $L_a$ there
with the pull-back of $L_a|_{D_a}$. In these terms $s_a=Az+B\bar z+O(|z|^2)$,
and we replace $s_a$ by $s_a-\chi B\bar z$, where $\chi$ is a cutoff
equal to $1$ near $D_a$ and supported in a small tubular neighbourhood.
During this change the section has the form
$Az+(1-\chi)B\bar z+O(|z|^2)$, whose first-order part has modulus at
least $(|A|-|B|)|z|$; so for a sufficiently small neighbourhood no new
zeros appear and transversality is preserved. The new normal derivative is $z\mapsto Az$, complex
linear. Perform this independently for all nodes; the isomorphisms
\eqref{eq:theta} are unaffected.

From this point on $\sigma$, $I$ and all the data
$(L_a,s_a,\theta_a,D_a)_{a\in\cT}$ are fixed.

\section{Recursive frames and their transition functions}\label{sec:frames}

Where $s_{a\pm}$ do not vanish, $F=\theta_a(s_{a+}\otimes s_{a-}^{-1})$
is a frame of $L_a$ and $f=s_a/F$ is a defining function of $D_a$;
nearer to $D_{a+}$ or $D_{a-}$ this frame degenerates and the small
child section must itself be replaced by the frame built from its own
children. The bookkeeping device is an \emph{active set}: the set of nodes
whose sections are ``small'' at the points under consideration, at
depth-dependent thresholds chosen so that a section is always much
smaller than the sections framing it.

\subsection*{Active sets and the multiscale cover}
Call a set $A$ of non-leaf nodes \emph{active} if it is
ancestor-closed: whenever a node belongs to $A$, all its ancestors do.
Its \emph{boundary} $\partial A$ consists of the roots not in $A$
together with the children of nodes of $A$ that do not belong to $A$
(these children may be leaves). The empty set is active, with
$\partial\varnothing=S$. There are finitely many active sets.

Choose Hermitian metrics on all $L_a$. We shall fix constants
$0<\eps_0<\eps_1<\dots<\eps_n$ in Section~\ref{sec:estimates}; they
will be small, with each ratio $\eps_d/\eps_{d+1}$ small. For an active
set $A$ define the open set
\begin{equation}\label{eq:UA}
 U_A=\Bigl\{x\in X:\
 |s_a(x)|<2\eps_{d(a)}\ \text{for }a\in A,\quad
 |s_b(x)|>\eps_{d(b)}\ \text{for }b\in\partial A\Bigr\}.
\end{equation}
Since $\{D_a\}_{a\in\cT}$ is a
positive SNC collection in dimension $2n$, no $n+1$ of the divisors
$D_a$ have a common point, and by compactness there is $\delta_0>0$
such that no $n+1$ distinct sections $s_a$ have norm $\le\delta_0$ at
the same point.

\begin{lemma}[The cover]\label{lem:cover}
If $2\eps_n<\delta_0$, then every nonempty $U_A$ has $|A|\le n$, and
the sets $U_A$, $A$ active, cover $X$.
\end{lemma}

\begin{proof}
If $x\in U_A$ then $|s_a(x)|<2\eps_{d(a)}\le2\eps_n<\delta_0$ for all
$a\in A$, so $|A|\le n$. For the covering, fix $x$ and determine $A$ by
the following process. Start at each root; call a node
\emph{accessible} if it is a root or the child of an activated node;
activate an accessible node $a$ precisely when $|s_a(x)|\le\eps_{d(a)}$,
and continue to its children only if it is activated. An accessible
leaf $\ell$ is never activated: otherwise $\ell$ and its $n$ ancestors
would be $n+1$ distinct nodes with sections of norm
$\le\eps_n<\delta_0$ at $x$. Hence the set $A$ of activated nodes is an
ancestor-closed set of non-leaf nodes, and its boundary consists
exactly of the accessible nodes that were not activated, at which
$|s_b(x)|>\eps_{d(b)}$. So $x\in U_A$.
\end{proof}

\subsection*{Frames and functions}
For an active set $A$ let
$O_A=\{x\in X: s_b(x)\ne0\ \text{for all }b\in\partial A\}$, an open
set containing $\overline{U_A}$, since $|s_b|\ge\eps_{d(b)}>0$ on
$\overline{U_A}$ for $b\in\partial A$. On $O_A$ define nonvanishing
sections $F^A_c$ of $L_c$, for $c\in A\cup\partial A$, recursively from
the boundary upwards, together with functions $f^A_a$ for $a\in A$:
\begin{align}
 F^A_b&=s_b &&(b\in\partial A),\label{eq:frames}\\
 F^A_a&=\theta_a\bigl(F^A_{a+}\otimes(F^A_{a-})^{-1}\bigr) &&(a\in A),\notag\\
 f^A_a&=s_a/F^A_a &&(a\in A).\label{eq:functions}
\end{align}
The recursion is well defined because for $a\in A$ each child of $a$
lies in $A\cup\partial A$, and the frames are nonvanishing by
induction. Each $f^A_a$ is an ordinary smooth complex-valued function
on $O_A$, and its zero set is $D_a\cap O_A$.

\begin{lemma}[Exact monomial transitions]\label{lem:monomial}
Let $A\subset B$ be active sets. Then for every $a\in A$ there are
integers $k_b$, $b\in\partial A\cap B$, such that
\begin{equation}\label{eq:monomial}
 f^A_a=f^B_a\prod_{b\in\partial A\cap B}(f^B_b)^{k_b}
 \qquad\text{on }O_A\cap O_B,
\end{equation}
where the functions $f^B_b$, $b\in\partial A\cap B$, are nonvanishing
on $O_A\cap O_B$, in particular on $U_A\cap O_B$.
\end{lemma}

\begin{proof}
First, every $b\in\partial A$ lies in $B\cup\partial B$: if $b\notin B$,
then $b$ is either a root not in $B$ or a child of a node of
$A\subset B$ not in $B$, hence $b\in\partial B$. Therefore $F^B_c$ is
defined for all $c\in A\cup\partial A$, and we may consider the ratios
\[
 r_c=F^A_c/F^B_c\in C^\infty(O_A\cap O_B,\C^*)
\qquad(c\in A\cup\partial A),
\]
quotients of two nonvanishing sections of the same line bundle $L_c$.
We show by induction from the boundary upwards that each $r_c$ is a
Laurent monomial with coefficient one in the $f^B_b$, $b\in\partial A\cap B$.
For $b\in\partial A$: if $b\notin B$ then $b\in\partial B$ and both
constructions use $s_b$ as the frame, so $r_b=1$; if $b\in B$ then
$r_b=s_b/F^B_b=f^B_b$, which is therefore nonvanishing on
$O_A\cap O_B$. For $a\in A$, complex linearity of $\theta_a$ gives
\[
 F^A_a=\theta_a\bigl(r_{a+}F^B_{a+}\otimes r_{a-}^{-1}(F^B_{a-})^{-1}\bigr)
 =\frac{r_{a+}}{r_{a-}}\,F^B_a,
 \qquad\text{i.e.}\qquad r_a=\frac{r_{a+}}{r_{a-}} ;
\]
the isomorphism $\theta_a$ cancels exactly. Finally
\[
 f^A_a=\frac{s_a}{F^A_a}=\frac{s_a}{F^B_a}\cdot\frac{F^B_a}{F^A_a}
 =f^B_a\,r_a^{-1},
\]
which is \eqref{eq:monomial}.
\end{proof}

By \eqref{eq:monomial}, for every almost complex structure $J$, if all
the functions $f^B_c$ occurring on the right-hand side are
$J$-holomorphic, then so is $f^A_a$; a transition $g\,f^B_a$ with an
arbitrary nonvanishing smooth $g$ would not have this property, which
is why actual bundle isomorphisms are used rather than unrelated local
trivializations.

\section{Uniform estimates and local tame structures}\label{sec:estimates}

The goal of this section is Proposition~\ref{prop:local}: near any
point where an active set $B$ is ``small and its boundary is not'',
there is a tame structure making all $f^B_a$, $a\in B$, holomorphic.
By Lemma~\ref{lem:graph} it suffices that the one-forms $F^B_a\,df^B_a$
have jointly surjective $(1,0)$-parts with uniformly bounded right
inverses and small $(0,1)$-parts. All constants are fixed \emph{before}
the thresholds $\eps_d$.

Fix Hermitian connections $\nabla_a$ on the bundles $L_a$ and an
$I$-Hermitian metric on $TX$. All norms below refer to these fixed
metrics. For an $L_a$-valued one-form $\alpha$ we write
$\alpha^{1,0}_I$ and $\alpha^{0,1}_I$ for its restrictions to
$T^{1,0}_IX$ and $T^{0,1}_IX$, so that $\alpha^{0,1}_I=0$ means that
$\alpha$ is $I$-complex linear.

\begin{lemma}[Uniform facts]\label{lem:uniform}
There are $\delta>0$ and $C_0,C_1>0$ such that:
\begin{enumerate}
\item No $n+1$ distinct sections $s_a$ have norm $\le\delta$ at a
common point.
\item Whenever $\Lambda\subset\cT$ and $x\in X$ satisfy $|s_a(x)|\le\delta$
for all $a\in\Lambda$, the complex-linear map
\begin{equation}\label{eq:PK}
 P_\Lambda=\bigl((\nabla_as_a)^{1,0}_I\bigr)_{a\in\Lambda}
 \colon T^{1,0}_{I,x}X\longrightarrow\bigoplus_{a\in\Lambda}(L_a)_x
\end{equation}
is surjective with a right inverse of norm at most $C_0$.
\item For every node $a$,
\begin{equation}\label{eq:anti}
 |(\nabla_as_a)^{0,1}_I|\le C_1|s_a|,\qquad|\nabla_as_a|\le C_1
 \quad\text{on }X.
\end{equation}
\end{enumerate}
\end{lemma}

\begin{proof}
(1) holds for every $\delta\le\delta_0$, by the choice of $\delta_0$
preceding Lemma~\ref{lem:cover}. For (2), fix
$\Lambda$ with $|\Lambda|\le n$ (by (1), larger $\Lambda$ do not occur).
On the compact common zero set $D_\Lambda=\bigcap_{a\in\Lambda}D_a$ the
connection derivatives $\nabla_as_a$ coincide with the intrinsic
derivatives $ds_a$, which are $I$-complex linear by
\eqref{eq:normal-linear} and jointly surjective onto
$\bigoplus_{a\in\Lambda}L_a$ by the transversality of the $D_a$; so
$P_\Lambda$ is a surjective complex-linear map there, with a right
inverse bounded by compactness. Surjectivity with a bound is an open
condition, so it persists on a neighbourhood of $D_\Lambda$, and by
compactness this neighbourhood contains
$\{|s_a|\le\delta\ \forall a\in\Lambda\}$ for small $\delta$; if
$D_\Lambda=\varnothing$, this set is empty for small $\delta$. There are
finitely many $\Lambda$. For (3), the second bound is compactness.
For the first, $(\nabla_as_a)^{0,1}_I$ vanishes on $D_a$ by
\eqref{eq:normal-linear}, so it is bounded by a constant times the
distance to $D_a$, and transversality of $s_a$ makes this distance
comparable to $|s_a|$ near $D_a$; away from $D_a$, $|s_a|$ is bounded
below.
\end{proof}

\subsection*{Differentials of the recursive frames}
Let $B$ be an active set and consider a point where
\begin{equation}\label{eq:closed-bounds}
 |s_a|\le2\eps_{d(a)}\quad(a\in B),\qquad
 |s_b|\ge\eps_{d(b)}\quad(b\in\partial B);
\end{equation}
these are the closed versions of the inequalities defining $U_B$, and
every point of $\overline{U_B}$ satisfies them. Such a point lies in
$O_B$, so the frames $F^B_c$ are defined near it. Their logarithmic
derivatives $\beta^B_c=(F^B_c)^{-1}\nabla_cF^B_c$, $c\in B\cup\partial B$,
are ordinary complex-valued one-forms. For each non-leaf node $a$, the
connection on $L_a$ transported by $\theta_a$ from the tensor-product
connection on $L_{a+}\otimes L_{a-}^{-1}$ differs from $\nabla_a$ by a
fixed smooth complex one-form $\gamma_a$, bounded on $X$. Hence the
recursion \eqref{eq:frames} gives
\begin{align*}
 \beta^B_a&=\beta^B_{a+}-\beta^B_{a-}+\gamma_a &&(a\in B),\\
 \beta^B_b&=s_b^{-1}\nabla_bs_b &&(b\in\partial B).
\end{align*}
Expanding the recursion, $\beta^B_a$ is a signed sum of boundedly many
forms $\gamma_c$ and of terminal terms $s_b^{-1}\nabla_bs_b$ with $b$ a
boundary descendant of $a$. If $a$ has depth $d$, all these boundary
descendants have depth at least $d+1$, so at a point satisfying
\eqref{eq:closed-bounds} the terminal terms are bounded by
$C_1/\eps_{d(b)}\le C_1/\eps_{d+1}$ by \eqref{eq:anti}. Assuming
$\eps_n\le1$, we obtain a constant $C_2$, depending only on $C_1$, the
forest and the $\gamma_c$, with
\begin{equation}\label{eq:beta}
 |\beta^B_a|\le C_2/\eps_{d(a)+1}\qquad(a\in B).
\end{equation}

Now multiply the differential of $f^B_a$ by its frame. Since
$s_a=f^B_aF^B_a$, the Leibniz rule gives
$\nabla_as_a=F^B_a\,df^B_a+f^B_aF^B_a\beta^B_a$, that is, the
$L_a$-valued one-form
\begin{equation}\label{eq:lambda}
 \lambda^B_a:=F^B_a\,df^B_a=\nabla_as_a-s_a\beta^B_a
\end{equation}
satisfies, by \eqref{eq:closed-bounds} and \eqref{eq:beta}, at points
satisfying \eqref{eq:closed-bounds},
\begin{equation}\label{eq:lambda-error}
 |\lambda^B_a-\nabla_as_a|\le2C_2\,\frac{\eps_{d(a)}}{\eps_{d(a)+1}}
 \qquad(a\in B),
\end{equation}
the size of $s_a$ divided by the size of the sections framing it.

\subsection*{Choice of the thresholds}
Choose
\begin{equation}\label{eq:thresholds}
 \eps_d=t^{\,n-d}\eps_n\qquad(0\le d\le n),\qquad 0<t<1,
\end{equation}
so that every ratio $\eps_d/\eps_{d+1}$ equals $t$, and take
$2\eps_n<\delta$ (hence also $2\eps_n<\delta_0$, as required in
Lemma~\ref{lem:cover}). By Lemma~\ref{lem:uniform}(1), every active set
$B$ for which \eqref{eq:closed-bounds} holds at some point has
$|B|\le n$. For such $B$ let $\lambda^B=(\lambda^B_a)_{a\in B}$, with
values in $\bigoplus_{a\in B}L_a$, and let
$P^\lambda_B=\lambda^B|_{T^{1,0}_IX}$,
$Q^\lambda_B=\lambda^B|_{T^{0,1}_IX}$. By \eqref{eq:lambda-error},
$P^\lambda_B$ differs from $P_B$ in \eqref{eq:PK} by at most
$2C_2\sqrt n\,t$, so for $t$ smaller than a constant depending on
$C_0,C_2,n$ it is surjective with a right inverse of norm at most
$2C_0$ (if $P'=P+\Delta$ with $\lVert\Delta\rVert\le1/(2C_0)$, then
$R(\id+\Delta R)^{-1}$ is a right inverse of $P'$ of norm $\le2C_0$).
By \eqref{eq:anti} and \eqref{eq:lambda-error},
\begin{equation}\label{eq:Qsmall}
 \lVert Q^\lambda_B\rVert\le\sqrt n\,(2C_1\eps_n+2C_2t)\le C_3(\eps_n+t).
\end{equation}
These bounds hold simultaneously for all such $B$ at all points
satisfying \eqref{eq:closed-bounds}. Now first choose $\eps_n$ small
and then $t$ small so that $2C_0C_3(\eps_n+t)<\mu_0$, where $\mu_0$ is
the constant of Lemma~\ref{lem:graph} for the fixed metrics. The
thresholds are fixed from now on.

\begin{proposition}[Local tame structures]\label{prop:local}
Let $B$ be an active set and $x$ a point satisfying
\eqref{eq:closed-bounds}. Then there are a neighbourhood $W\subset O_B$
of $x$ and a smooth $\sigma$-tame almost complex structure on $W$ for
which all the functions $f^B_a$, $a\in B$, are holomorphic.
\end{proposition}

\begin{proof}
Apply Lemma~\ref{lem:graph} with $V=\bigoplus_{a\in B}L_a$,
$\lambda=\lambda^B$, $P=P^\lambda_B$, $Q=Q^\lambda_B$ and
$R=P^*(PP^*)^{-1}$. At $x$ we have $\lVert R\rVert\le2C_0$ (some
right inverse has norm $\le2C_0$, and $R$ has the least norm) and
$\lVert Q\rVert\le C_3(\eps_n+t)$, so $\lVert\mu\rVert<\mu_0$.
Surjectivity of $P$ and the bound on $\mu$ are open conditions, so
they hold on a neighbourhood $W$ of $x$, and the lemma provides
$J_\mu$ on $W$; since $\lambda^B_a=F^B_a\,df^B_a$ with nonvanishing
$F^B_a$, each $f^B_a$ is $J_\mu$-holomorphic.
\end{proof}

\section{Simultaneous local choices and global patching}\label{sec:patching}

To patch the local structures of Proposition~\ref{prop:local} we must
make sure that near a point $x$ \emph{all} the functions $f^A_a$ of
\emph{all} the sets $U_A$ that come close to $x$ are holomorphic for
the local structure chosen at $x$. Lemma~\ref{lem:monomial} reduces
this to the functions of one maximal active set. For $x\in X$ let
\begin{equation}\label{eq:Bx}
 B_x=\bigcup\{A:\ A\text{ active},\ x\in\overline{U_A}\}.
\end{equation}
A union of ancestor-closed sets of non-leaf nodes is again one, so
$B_x$ is an active set ($B_x=\varnothing$ is allowed).

\begin{lemma}\label{lem:Bx}
The point $x$ satisfies the inequalities \eqref{eq:closed-bounds} for
$B=B_x$. In particular $|B_x|\le n$ and the frames $F^{B_x}_c$ are
defined near $x$.
\end{lemma}

\begin{proof}
If $a\in B_x$, then $a\in A$ for some $A$ with $x\in\overline{U_A}$,
so $|s_a(x)|\le2\eps_{d(a)}$. For the boundary inequalities, let
$b\in\partial B_x$. If $b$ is a child of some $a\in B_x$ with
$b\notin B_x$, choose $A$ with $a\in A$ and $x\in\overline{U_A}$; then
$b\notin A$ because $A\subset B_x$, so $b\in\partial A$ and
$|s_b(x)|\ge\eps_{d(b)}$. If $b$ is a root not in $B_x$, then $b$ is a
root not in $A$ for every $A$ occurring in \eqref{eq:Bx}, so
$b\in\partial A$ and again $|s_b(x)|\ge\eps_{d(b)}$. The remaining
assertions follow from Lemma~\ref{lem:uniform}(1) and from
$2\eps_n<\delta$, as in Section~\ref{sec:estimates}.
\end{proof}

By Lemma~\ref{lem:Bx} and Proposition~\ref{prop:local} there are a
neighbourhood $W_x\subset O_{B_x}$ of $x$ and a $\sigma$-tame structure
$J_x$ on $W_x$ for which every $f^{B_x}_a$, $a\in B_x$, is holomorphic
(for $B_x=\varnothing$ take $J_x=I$). Shrink $W_x$ so that
\begin{equation}\label{eq:shrink}
 U_A\cap W_x\ne\varnothing\quad\Longrightarrow\quad A\subset B_x .
\end{equation}
This is possible: there are finitely many active sets, and if
$x\notin\overline{U_A}$ then some neighbourhood of $x$ misses $U_A$,
while $x\in\overline{U_A}$ gives $A\subset B_x$ by definition.

For every active $A$ with $U_A\cap W_x\ne\varnothing$ we thus have
$A\subset B_x$, and Lemma~\ref{lem:monomial} expresses each $f^A_a$ on
$U_A\cap W_x\subset U_A\cap O_{B_x}$ as a Laurent monomial, with
nonvanishing denominators, in the $J_x$-holomorphic functions
$f^{B_x}_c$. Hence
\begin{equation}\label{eq:all-local}
 f^A_a\ \text{is $J_x$-holomorphic on }U_A\cap W_x
 \qquad(A\text{ active},\ a\in A).
\end{equation}

\begin{proposition}[The global structure]\label{prop:global}
There is a $\sigma$-tame almost complex structure $J$ on $X$ such that
\begin{equation}\label{eq:finalhol}
 df^A_a\circ J=i\,df^A_a\quad\text{on }U_A
 \qquad\text{for every active set $A$ and every }a\in A.
\end{equation}
\end{proposition}

\begin{proof}
Choose finitely many points $x_1,\dots,x_N$ with
$X=\bigcup_kW_{x_k}$, a partition of unity $(\rho_k)$ subordinate to
this cover, and a $\sigma$-compatible reference structure $K$ (for
instance $K=I$, which is compatible). Let $J$ be the tame structure of
Lemma~\ref{lem:convex} with graph map
$\Phi_J=\sum_k\rho_k\Phi_{J_{x_k}}$. Let $y\in U_A$ and $a\in A$. Every $k$
with $\rho_k(y)\ne0$ has $y\in U_A\cap W_{x_k}$, so $f^A_a$ is
$J_{x_k}$-holomorphic near $y$ by \eqref{eq:all-local}. By
Lemma~\ref{lem:convex}, $f^A_a$ is $J$-holomorphic at $y$.
\end{proof}

\section{The coordinates and the deformation}\label{sec:end}

\begin{proof}[Proof of Theorem~\ref{thm:main}]
\emph{The coordinates.} Let $p\in D$ and let $D_{i_1},\dots,D_{i_r}$
be the original branches through $p$; $1\le r\le n$. By
Lemma~\ref{lem:cover} choose an active set $A$ with $p\in U_A$. Each
root $i_j$ belongs to $A$: a root not in $A$ lies in $\partial A$ and
its section is nonzero on $U_A$, whereas $s_{i_j}(p)=0$. Define
$z_j=f^A_{i_j}$ for $1\le j\le r$: smooth functions on $U_A$,
$J$-holomorphic by \eqref{eq:finalhol}, with $z_j^{-1}(0)=D_{i_j}\cap U_A$
because the frame $F^A_{i_j}$ does not vanish. At $p$, since
$s_{i_j}(p)=0$, $dz_j(p)=(F^A_{i_j}(p))^{-1}ds_{i_j}(p)$, and the real
derivatives $ds_{i_1}(p),\dots,ds_{i_r}(p)$ are jointly surjective onto
$\bigoplus_j(L_{i_j})_p$ by the transversality of the branches through
$p$ (Definition~\ref{def:snc}). So $(z_1,\dots,z_r)\colon U_A\to\C^r$
is a submersion at $p$, and the implicit function theorem completes it
by smooth functions $z_{r+1},\dots,z_n$ to a smooth complex coordinate
system on a neighbourhood of $p$. Shrink this
neighbourhood so that it misses the branches not containing $p$. On it,
$D=D_{i_1}\cup\dots\cup D_{i_r}=\{z_1\cdots z_r=0\}$. For $p\notin D$
take a neighbourhood disjoint from $D$, arbitrary smooth complex
coordinates and $r=0$.

\emph{The deformation.} The only changes of the symplectic form in the
whole construction were the regularization deformations of
Lemma~\ref{lem:reg}, applied in Section~\ref{sec:forest}: one inside
each of the $m(2^n-1)$ applications of Lemma~\ref{lem:aux} and one at
the end. Each is a path $\sigma+d\tau_t$ of forms cohomologous to its
starting point along which every divisor constructed up to that stage,
in particular every original branch $D_i$, remains part of a positive
symplectic SNC collection. Concatenating these finitely many paths,
reparametrized smoothly at the joins, gives the path $\omega_t$ of
Theorem~\ref{thm:main}: $\omega_0=\omega$, $[\omega_t]=[\omega]$,
$\{D_i\}$ is positive SNC for every $\omega_t$, and $\omega_1=\sigma$.
The structure $J$ of Proposition~\ref{prop:global} is $\sigma$-tame,
that is, $\omega_1$-tame. This completes the proof.
\end{proof}

\begin{thebibliography}{99}

\bibitem{Auroux1997}
D.~Auroux,
\emph{Asymptotically holomorphic families of symplectic submanifolds},
Geom.\ Funct.\ Anal.\ \textbf{7} (1997), 971--995.

\bibitem{Auroux2001b}
D.~Auroux,
\emph{Symplectic maps to projective spaces and symplectic invariants},
Turkish J.\ Math.\ \textbf{25} (2001), 1--42; arXiv:math/0007130.

\bibitem{Auroux2001}
D.~Auroux,
\emph{Estimated transversality in symplectic geometry and projective maps},
in: Symplectic Geometry and Mirror Symmetry (Seoul, 2000), World
Scientific, 2001, 1--30; arXiv:math/0010052.

\bibitem{CM}
K.~Cieliebak and K.~Mohnke,
\emph{Symplectic hypersurfaces and transversality in Gromov--Witten theory},
J.\ Symplectic Geom.\ \textbf{5} (2007), 281--356.

\bibitem{Donaldson1996}
S.~K.~Donaldson,
\emph{Symplectic submanifolds and almost-complex geometry},
J.\ Differential Geom.\ \textbf{44} (1996), 666--705.

\bibitem{FMZ}
M.~Farajzadeh Tehrani, M.~McLean, and A.~Zinger,
\emph{Normal crossings singularities for symplectic topology},
Adv.\ Math.\ \textbf{339} (2018), 672--748; arXiv:1410.0609v3.

\bibitem{FMZS}
M.~Farajzadeh Tehrani, M.~McLean, and A.~Zinger,
\emph{Normal crossings singularities for symplectic topology: structures},
Acta Math.\ Sin.\ (Engl.\ Ser.) \textbf{40} (2024), 107--160;
arXiv:2112.13125.

\bibitem{FZ}
M.~Farajzadeh Tehrani and A.~Zinger,
\emph{On symplectic sum formulas in Gromov--Witten theory},
arXiv:1404.1898 (2014).

\bibitem{Mohsen2019}
J.-P.~Mohsen,
\emph{Transversalit\'e quantitative en g\'eom\'etrie symplectique:
sous-vari\'et\'es et hypersurfaces},
Ann.\ Fac.\ Sci.\ Toulouse Math.\ (6) \textbf{28} (2019), no.~4, 655--706.

\bibitem{Wendl}
C.~Wendl,
\emph{Lectures on holomorphic curves in symplectic and contact geometry},
arXiv:1011.1690v2 (2014).

\end{thebibliography}
\end{document}
